\section{Байесовы игры: приложения}

\subsection{Выборы из двух кандидатов с издержками голосования}

Рассмотрим однопериодную игру, моделирующую процесс голосования. В выборах участвуют $N_1 + N_2 = N$ избирателей, где $N_1$ — сторонники кандидата $A$, а $N_2$ — сторонники кандидата $B$. 

Каждый избиратель принимает решение: проголосовать за одного из двух кандидатов ($A$ или $B$) либо остаться дома (стратегия $o$). Множество стратегий избирателя имеет вид $S_i = \{A, B, o\}$. 

Выигрыш избирателя формируется следующим образом:
\begin{itemize}
    \item В случае победы предпочитаемого кандидата избиратель получает полезность, равную $1$.
    \item В случае победы оппонента избиратель получает $0$.
    \item Участие в голосовании сопряжено с индивидуальными издержками $c_i < \frac{1}{2}$.
\end{itemize}

Вероятность победы кандидата $A$ определяется правилом большинства:
\begin{equation}
P_A = 
\begin{cases} 
1, & \#\{i \mid s_i = A\} > \#\{i \mid s_i = B\} \\ 
\frac{1}{2}, & \#\{i \mid s_i = A\} = \#\{i \mid s_i = B\} \\ 
0, & \#\{i \mid s_i = A\} < \#\{i \mid s_i = B\} 
\end{cases}
\end{equation}
где $\#$ обозначает мощность множества (количество голосов). Вероятность победы второго кандидата составляет $P_B = 1 - P_A$. Заметим, что голосование за непредпочтительного кандидата является строго доминируемой стратегией.

\subsubsection{Случай полной информации}

Предположим, что издержки всех избирателей одинаковы и общеизвестны ($c_i = c$), а размеры групп равны ($N_1 = N_2 = N$). Будем искать симметричное равновесие в смешанных стратегиях, в котором каждый избиратель принимает участие в голосовании с вероятностью $q$.

Для сторонника кандидата $A$ условие безразличия между чистыми стратегиями (голосовать и остаться дома) в равновесии принимает вид:
\begin{equation}
P_A(1, q) - c = P_A(0, q) \implies c = P_A(1, q) - P_A(0, q)
\end{equation}
где $P_A(1, q)$ — вероятность победы $A$ при условии, что данный избиратель проголосовал, а $P_A(0, q)$ — вероятность победы $A$, если данный избиратель остался дома. 

Вычисление данных вероятностей через биномиальные коэффициенты дает:
\begin{align}
P_A(1, q) &= \sum_{i=0}^{N-1} \left( C_{N-1}^i q^i (1 - q)^{N-i-1} \sum_{j=0}^{i} C_N^j q^j (1 - q)^{N-j} \right) + \nonumber \\
&+ \frac{1}{2} \sum_{i=0}^{N-1} \left( C_{N-1}^i C_N^{i+1} q^{2i+1} (1 - q)^{2N-2i-2} \right)
\end{align}

\begin{align}
P_A(0, q) &= \sum_{i=1}^{N-1} \left( C_{N-1}^i q^i (1 - q)^{N-i-1} \sum_{j=0}^{i-1} C_N^j q^j (1 - q)^{N-j} \right) + \nonumber \\
&+ \frac{1}{2} \sum_{i=0}^{N-1} \left( C_{N-1}^i C_N^i q^{2i} (1 - q)^{2N-2i-1} \right)
\end{align}

Анализ данного уравнения показывает, что для любого $c < \frac{1}{2}$ при достаточно большом $N$ существует два решения:
\begin{enumerate}
    \item Равновесие с низкой явкой: при $N \to \infty \implies q \to 0$.
    \item Равновесие с высокой явкой: при $N \to \infty \implies q \to 1$.
\end{enumerate}
При асимметрии групп ($N_A \neq N_B$) теоретические модели предсказывают, что в группе меньшего размера явка будет выше.

\subsubsection{Случай неполной информации (частная информация об издержках)}

Пусть индивидуальные издержки $c_i$ являются случайной величиной, распределенной на отрезке $[0, 1]$ с кумулятивной функцией распределения $F(\cdot)$. Стратегия избирателя теперь представляет собой функцию от его издержек. 

Рассмотрим симметричную стратегию. Пусть все остальные избиратели приходят на выборы с вероятностью $q$. Отдельный избиратель с издержками $c$ пойдет на выборы тогда и только тогда, когда ожидаемая выгода превышает издержки:
\begin{equation}
P_A(1, q) - c > P_A(0, q) \implies c < \bar{c} = P_A(1, q) - P_A(0, q)
\end{equation}
Это монотонная чистая стратегия: избиратель голосует, если его издержки ниже порогового значения $\bar{c}$. В силу симметрии, вероятность того, что случайно выбранный избиратель проголосует, равна:
\begin{equation}
q = P(c < P_A(1, q) - P_A(0, q)) = F(P_A(1, q) - P_A(0, q))
\end{equation}
Откуда следует условие равновесия: $F^{-1}(q) = P_A(1, q) - P_A(0, q)$. В данной модели вероятность участия (явка) монотонно убывает с ростом размера популяции $N$.

\begin{example}[Экспериментальное подтверждение]
В работе Levine and Palfrey (2007) проводилось лабораторное тестирование описанной модели. Издержки генерировались из равномерного распределения на интервале $[0, 1/2]$. 

Теоретические предсказания:
\begin{itemize}
    \item Явка внутри группы уменьшается с ростом ее размера.
    \item Явка уменьшается с ростом стоимости голосования.
    \item Явка меньше в большей группе (эффект размера большинства).
\end{itemize}

Результаты сравнения теоретических вероятностей участия ($p^*$) и эмпирических частот ($\widehat{p}$) представлены в Таблице \ref{tab:voting_experiment}. Данные подтверждают основные качественные предсказания модели.

\begin{table}[h]
\centering
\renewcommand{\arraystretch}{1.2}
\begin{NiceTabular}{cc|c|ccc|ccc}
\toprule
$N$ & $N_A$ & $N_B$ & $\widehat{p}_A$ & $p_A^*$ & $p_A^{\lambda=7}$ & $\widehat{p}_B$ & $p_B^*$ & $p_B^{\lambda=7}$ \\
\midrule
$3$ & $1$ & $2$ & $0.530$ $(0.017)$ & $0.537$ & $0.549$ & $0.593$ $(0.012)$ & $0.640$ & $0.616$ \\
$9$ & $3$ & $6$ & $0.436$ $(0.013)$ & $0.413$ & $0.421$ & $0.398$ $(0.009)$ & $0.374$ & $0.395$ \\
$9$ & $4$ & $5$ & $0.479$ $(0.012)$ & $0.460$ & $0.468$ & $0.451$ $(0.010)$ & $0.452$ & $0.463$ \\
$27$ & $9$ & $18$ & $0.377$ $(0.011)$ & $0.270$ & $0.297$ & $0.282$ $(0.007)$ & $0.228$ & $0.275$ \\
$27$ & $13$ & $14$ & $0.385$ $(0.009)$ & $0.302$ & $0.348$ & $0.356$ $(0.009)$ & $0.297$ & $0.345$ \\
$51$ & $17$ & $34$ & $0.333$ $(0.011)$ & $0.206$ & $0.245$ & $0.266$ $(0.008)$ & $0.171$ & $0.230$ \\
$51$ & $25$ & $26$ & $0.390$ $(0.010)$ & $0.238$ & $0.301$ & $0.362$ $(0.009)$ & $0.235$ & $0.300$ \\
\bottomrule
\end{NiceTabular}
\caption{Сравнение теоретической модели и экспериментальных данных (Levine, Palfrey, 2007)}
\label{tab:voting_experiment}
\end{table}
\end{example}

\subsection{Приближение смешанных равновесий чистыми}

\begin{theorem}[Очищение по Харшаньи]
Дж. Харшаньи (1973) доказал, что практически в любой игре с полной информацией равновесие в смешанных стратегиях можно рассматривать как предельный случай равновесия в чистых стратегиях в некоторой последовательности игр с неполной информацией, когда дисперсия частных шоков выигрышей стремится к нулю.
\end{theorem}

\begin{example}
Рассмотрим игру $2 \times 2$ с полной информацией, матрица выигрышей которой имеет следующий вид:
\begin{table}[h]
\centering
\renewcommand{\arraystretch}{1.3}
\begin{NiceTabular}{c|cc}
\toprule
\textbf{Игрок 1 \textbackslash\ Игрок 2} & \textbf{A} & \textbf{B} \\
\midrule
\textbf{A} & $2, 2$ & $0, 0$ \\
\textbf{B} & $0, 0$ & $1, 2$ \\
\bottomrule
\end{NiceTabular}
\end{table}

Введем неполную информацию. Пусть $\epsilon > 0$. Выигрыш игрока $i$ зависит от его собственного частного типа $t_i$, где величины $t_1, t_2$ независимо и равномерно распределены на отрезке $[-1, 1]$. Модифицированная функция выигрыша задается как:
\begin{equation}
\tilde{u}_i(a_i, a_{-i}) = 
\begin{cases}
u_i(A, a_{-i}) + \epsilon t_i, & a_i = A \\
u_i(B, a_{-i}), & a_i = B
\end{cases}
\end{equation}

Пусть $p_1, p_2$ — вероятности того, что игроки 1 и 2 соответственно выберут действие $A$. При заданном $p_2$, игрок 1 предпочтет выбрать $A$ тогда и только тогда:
\begin{equation}
p_2 \cdot (2 + \epsilon t_1) + (1 - p_2) \epsilon t_1 \ge p_2 \cdot 0 + (1 - p_2) \cdot 1
\end{equation}
После упрощения получаем условие на тип $t_1$:
\begin{equation}
\epsilon t_1 \ge 1 - 3p_2
\end{equation}
Аналогичное рассуждение для второго игрока дает:
\begin{equation}
p_1 \cdot (2 + \epsilon t_2) + (1 - p_1) \epsilon t_2 \ge p_1 \cdot 0 + (1 - p_1) \cdot 2 \implies \epsilon t_2 \ge 2 - 4p_1
\end{equation}

Опираясь на равномерное распределение $t_1, t_2$, можно построить функции наилучшего ответа (вероятности выбора $A$):
\begin{equation}
\check{p}_1(p_2) = P(\epsilon t_1 \ge 1 - 3p_2) = 
\begin{cases}
0, & p_2 \le \frac{1-\epsilon}{3} \\
\frac{1}{2} - \frac{1-3p_2}{2\epsilon}, & p_2 \in \left(\frac{1-\epsilon}{3}, \frac{1+\epsilon}{3}\right) \\
1, & p_2 \ge \frac{1+\epsilon}{3}
\end{cases}
\end{equation}

\begin{equation}
\check{p}_2(p_1) = P(\epsilon t_2 \ge 2 - 4p_1) = 
\begin{cases}
0, & p_1 \le \frac{2-\epsilon}{4} \\
\frac{1}{2} - \frac{1-2p_1}{\epsilon}, & p_1 \in \left(\frac{2-\epsilon}{4}, \frac{2+\epsilon}{4}\right) \\
1, & p_1 \ge \frac{2+\epsilon}{4}
\end{cases}
\end{equation}

При $\epsilon < \sqrt{2}$ данная система уравнений имеет три решения:
\begin{enumerate}
    \item $(p_1^*, p_2^*) = (0, 0)$ — оба игрока играют $(B, B)$.
    \item $(p_1^*, p_2^*) = (1, 1)$ — оба играют $(A, A)$.
    \item Внутреннее решение, задаваемое системой нелинейных уравнений:
\end{enumerate}
\begin{equation}
p_1^* = \frac{\frac{3}{2} + \frac{\epsilon}{4} - \frac{\epsilon^2}{2}}{3 - \epsilon^2}, \quad p_2^* = \frac{1 - \frac{\epsilon^2}{2}}{3 - \epsilon^2}
\end{equation}

Заметим, что при устремлении дисперсии шоков к нулю ($\epsilon \to 0$), внутреннее решение сходится к значениям:
\begin{equation}
(p_1^*, p_2^*) \to \left(\frac{1}{2}, \frac{1}{3}\right)
\end{equation}
что в точности совпадает с равновесием в смешанных стратегиях в исходной игре с полной информацией.
\end{example}

\subsection{Обработка сигналов в правосудии (Суд присяжных)}

Рассмотрим модель принятия решений коллегиальным органом. Двое судей должны вынести вердикт подсудимому: оправдать ($I$) или обвинить ($G$).

Состояние мира $t$ скрыто от судей. Подсудимый в действительности либо совершил преступление ($t = G$), либо невиновен ($t = I$). Априорная вероятность невиновности равна $P(t = I) = p$.

Каждый судья $i \in \{1, 2\}$ получает частный независимый сигнал $\theta_i \in \{I, G\}$ о виновности. Структура сигналов задается следующим образом:
\begin{itemize}
    \item Если $t = I$, сигнал безошибочен: $\theta_i = I$ с вероятностью $1$.
    \item Если $t = G$, сигнал зашумлен: $\theta_i = I$ с вероятностью $\frac{1}{2}$ и $\theta_i = G$ с вероятностью $\frac{1}{2}$.
\end{itemize}

Процедура принятия решения такова: каждый судья выносит свой вердикт $a_i \in \{I, G\}$ независимо и одновременно. Для осуждения требуется единогласие: подсудимый признается виновным, только если $a_1 = a_2 = G$. Выигрыш каждого судьи равен $1$ при правильном решении суда и $0$ в противном случае. 

Рассмотрим множество состояний мира и их совместные вероятности, учитывая структуру сигналов:
\begin{table}[h]
\centering
\renewcommand{\arraystretch}{1.5}
\begin{NiceTabular}{l|cc}
\toprule
 & $\theta_2 = G$ & $\theta_2 = I$ \\
\midrule
$\theta_1 = I, t = G$ & $\frac{1-p}{4}$ & $\frac{1-p}{4}$ \\
$\theta_1 = G, t = G$ & $\frac{1-p}{4}$ & $\frac{1-p}{4}$ \\
$\theta_1 = I, t = I$ & $0$ & $p$ \\
$\theta_1 = G, t = I$ & $0$ & $0$ \\
\bottomrule
\end{NiceTabular}
\end{table}

Используя формулу Байеса, обновим убеждения судьи $i$ после получения сигнала $\theta_i = I$:
\begin{equation}
P(t = I \mid \theta_i = I) = \frac{p}{\frac{1-p}{4} + \frac{1-p}{4} + p} = \frac{2p}{1 + p}
\end{equation}
Если же судья получает сигнал $\theta_i = G$, то он абсолютно уверен в виновности: $P(t = I \mid \theta_i = G) = 0$. Следовательно, выбор $a(G) = G$ является слабо доминируемой стратегией при таком сигнале.

Остается определить оптимальное поведение при получении сигнала $\theta_i = I$. Пусть второй судья при сигнале $\theta_2 = I$ голосует за оправдание с вероятностью $q_2 = P(a_2 = I)$. Оценим ожидаемые полезности первого судьи от выбора стратегий $I$ и $G$, полагая, что сигнал $\theta_1 = I$:
\begin{equation}
Eu_1(I, q_2 \mid \theta_1 = I) = P(t = I, a_2 = I \mid \theta_1 = I) + P(t = I, a_2 = G \mid \theta_1 = I) = \frac{2p}{1 + p}
\end{equation}
\begin{equation}
Eu_1(G, q_2 \mid \theta_1 = I) = \frac{2p}{1 + p} q_2 + \frac{1 - p}{1 + p} \frac{2 - q_2}{2}
\end{equation}

Условие, при котором оправдание строго предпочтительнее обвинения ($u_1(I) > u_1(G)$), сводится к неравенству:
\begin{equation}
q_2(1 - 5p) > 2 - 6p
\end{equation}
Анализ данного неравенства позволяет построить функцию наилучшего ответа $\check{q}_1(q_2)$. При $p \ge \frac{1}{3}$:
\begin{equation}
\check{q}_1(q_2) = 
\begin{cases}
0, & q_2 > \frac{6p - 2}{5p - 1} \\
[0, 1], & q_2 = \frac{6p - 2}{5p - 1} \\
1, & q_2 < \frac{6p - 2}{5p - 1}
\end{cases}
\end{equation}
Симметричное равновесие в данном случае достигается при $q_1^* = q_2^* = \frac{6p - 2}{5p - 1}$. 

Если же априорная вероятность невиновности мала ($p < \frac{1}{3}$), доминирующей стратегией при любом сигнале становится обвинение ($q_1 = 0$). В таком равновесии информативные сигналы игнорируются, и суд всегда выносит обвинительный вердикт.

\textbf{Эффективность равновесия.} Вычислим вероятности судебных ошибок:
\begin{itemize}
    \item Осуждение невиновного (ошибка I рода):
    \begin{equation}
    P(a_1 = G, a_2 = G \mid t = I) = 
    \begin{cases}
    1, & p \le \frac{1}{3} \\
    \frac{(1-p)^2}{(5p-1)^2}, & p > \frac{1}{3}
    \end{cases}
    \end{equation}
    \item Оправдание виновного (ошибка II рода):
    \begin{equation}
    P(a_1 = I \text{ или } a_2 = I \mid t = G) = \frac{q^* - (q^*)^2}{2}
    \end{equation}
\end{itemize}
